%! TeX root = ../main.tex

\chapter{Conclusion and outlook}%
\label{ch:conclusion}

In this thesis we developed and implemented a DMFT impurity solver
which works on the real axis at zero temperature.
By utilizing the natural impurity orbitals
and restricted active space approach~\cite{Lu2014, Lu2019},
we successfully transform our problem from an exponentially scaling one to a polynomial one.
While previously a few number of bath sites was the limiting factor for ED,
this approach allows us to couple the impurity to hundreds of bath sites.
To calculate the self-energy, we use the \emph{symmetric improved estimator}~\cite{Kugler2022}
without any broadening using the Schur complement.

We then tested the code on the half-filled Bethe lattice in infinite dimensions.
By continuously increasing the Coulomb interaction $U$, we obtain a critical value
$U_{c2}/D \in [3.0, 3.3]$ for the Mott metal-to-insulator transition.
The exact value depends on the discretization of the hybridization function and method
used to extract $U_{c2}$.
For intermediate interaction strength,
the pinning condition $\muppi DA(\omega=0)=2.0$ is well fulfilled.
We observe small peaks at the inner edge
of the Hubbard bands, comparable to previous results~\cite{Karski2005, Karski2008, Ganahl2015}.
To obtain the quasi-particle weight $Z$, we can directly
use the derivative of the real part of the self-energy
without taking a finite difference quotient.
Comparison of the low order moments of the Green's function and self-energy
show excellent agreement against analytical results.

The solver was then briefly tested on a one-orbital model of CeRu$_4$Sn$_6$.
As a finite broadening was used, the height of the Kondo peak was smaller than anticipated.
It may be possible to improve this calculation by using the approach based on the Schur complement.

There are multiple directions to improve and enhance the current code:
First,
it was not possible to find a metastable insulating solution.
A variety of approaches were unsuccessfully tested: mixing of the self-energy between steps
with a parameter $\alpha$ as low as \num{0.01};
using a Hubbard-III approximation for the initial hybridization function;
starting with an insulating self-energy with all weight concentrated at a single pole
at the Fermi-level.

Second, multithreading performance can be improved.
Currently, using lots of cores improves the speed-up by a factor of less than two.
Applying the Hamiltonian onto a wave function will linearly combine Slater determinants
leading to memory conflict.
Here, one must find a performant algorithm without race conditions.

Last, for realistic material calculations one needs to expand the code to multi-orbital systems.
As this will further increase computational complexity,
optimizations mentioned in the previous step should be done first.

The code discussed in this thesis is contained in two packages:
\texttt{Fermions.jl}~\cite{Wallerberger2022}
(not yet public)
and \texttt{RAS\_DMFT.jl} (\url{https://github.com/frankebel/RAS_DMFT.jl}).
